% --- MODULE 2: Linear Data Structures ---
% --- LECTURES 4, 5, 6, 7, 8, 10 & CLRS 10 ---

\section{Module 2: Linear Data Structures (Lists, Stacks, Queues)}

\subsection{Arrays vs. Linked Lists (Lec 4, 5, 6)}
A fundamental trade-off in data structures.
\begin{center}
\begin{tabular}{|l|l|l|}
\hline
\textbf{Operation} & \textbf{Array} & \textbf{Linked List} \\ \hline
Read (by index) & $O(1)$ & $O(n)$ \\
Search & $O(n)$ & $O(n)$ \\
Insert (at front) & $O(n)$ & $O(1)$ \\
Insert (at end) & $O(1)$* & $O(1)$** \\
Delete (at front) & $O(n)$ & $O(1)$ \\
Delete (at end) & $O(1)$ & $O(n)$*** \\
\hline
\end{tabular}
\end{center}
\begin{footnotesize}
*Amortized $O(1)$ with dynamic array (Lec 9).
**Requires a \texttt{tail} pointer.
***Requires traversal to find predecessor, or a Doubly Linked List.
\end{footnotesize}
\begin{itemize}
    \item \textbf{Sentinel Nodes (CLRS 10.2):} To simplify code, use a special dummy node `L.nil` for the `head` and to mark the `next` of the tail. This removes edge cases for empty lists or single-node lists.
\end{itemize}

\subsection{Stack ADT (LIFO - Last-In, First-Out) (Lec 6, 7)}
\begin{description}
    \item[Push(e)] Insert element at the top. \textbf{$O(1)$*}.
    \item[Pop()] Remove and return element from top. \textbf{$O(1)$}.
    \item[Top()] Return top element without removing. \textbf{$O(1)$}.
\end{description}
\begin{footnotesize}
*Amortized $O(1)$ using a dynamic array (Lec 9).
\end{footnotesize}

\subsubsection{Application: Balanced Parentheses (Lec 7)}
Algorithm to check if a string of brackets is balanced.
\begin{itemize}
    \item Create an empty stack.
    \item Iterate through the string:
    \begin{itemize}
        \item If char is an \textbf{opening} bracket ('(', '[', '\{'), \textbf{push} it onto the stack.
        \item If char is a \textbf{closing} bracket (')', ']', '\}'):
        \begin{itemize}
            \item If stack is empty $\to$ \textbf{return false}.
            \item \textbf{Pop} the top element.
            \item If popped bracket does not match the closing bracket $\to$ \textbf{return false}.
        \end{itemize}
    \end{itemize}
    \item After the loop, \textbf{return true} if the stack is empty, else \textbf{false}.
\end{itemize}
\textbf{Complexity:} $O(n)$ time, $O(n)$ space.

\subsection{Queue ADT (FIFO - First-In, First-Out) (Lec 8)}
\begin{description}
    \item[Enqueue(e)] Insert element at the back. \textbf{$O(1)$*}.
    \item[Dequeue()] Remove and return element from front. \textbf{$O(1)$**}.
    \item[Front()] Return front element without removing. \textbf{$O(1)$}.
\end{description}
\begin{footnotesize}
*Amortized $O(1)$ using a dynamic/doubling array.
**Requires a \textbf{circular array} implementation to avoid an $O(n)$ shift.
\end{footnotesize}

\subsubsection{Application: Maze Solving / BFS (Lec 10)}
\begin{itemize}
    \item Queues are the core data structure for \textbf{Breadth-First Search (BFS)}.
    \item BFS explores "layer by layer", which is used to find the \textbf{shortest path} in unweighted graphs (like a maze).
    \item \textbf{Proof of Correctness:} This is a non-trivial proof by induction. See the proof in \textbf{Module 7} (Graphs).
\end{itemize}

\subsection{General Tree Representation (CLRS 10.3)}
How to store a tree where nodes have an \textit{arbitrary} number of children?
\begin{itemize}
    \item \textbf{Bad Idea:} Array of pointers at each node (wastes space).
    \item \textbf{Good Idea: Left-Child, Right-Sibling Representation.}
    \item Each node $x$ has only two pointers:
    \begin{enumerate}
        \item $x.\text{left\_child}$: points to its first child.
        \item $x.\text{right\_sibling}$: points to its next sibling.
    \end{enumerate}
    \item This creates a binary tree, where all "left" pointers go down one level, and all "right" pointers stay on the same level.
\end{itemize}